\documentclass[12pt,a4paper]{article}
\usepackage{geometry}
\geometry{margin=2.54cm}
\usepackage{graphicx}
\usepackage{booktabs}
\usepackage{grffile}
\usepackage{xcolor}
\usepackage{hyperref}
\hypersetup{colorlinks=true, linkcolor=blue, urlcolor=blue}
\usepackage{xeCJK}
\setCJKmainfont{SimSun}
\setCJKsansfont{SimHei}
\setCJKmonofont{FangSong}
\title{600580 卧龙电驱 多智能体研究报告}
\author{FinAgent}
\date{\today}
\begin{document}
\maketitle
\tableofcontents
\newpage
\sloppy
\section{600580 卧龙电驱 多智能体研究报告}

> **草稿状态**：验证评分 85，结论 `pass\_after\_revision`。 本报告尚未通过验证 Agent 复核，请勿对外发布。

\subsection{ 目录}

\begin{itemize}
\item [执行摘要](\#执行摘要)
\item [核心指标速览](\#核心指标速览)
\item [量价与技术面](\#量价与技术面)
\item [基本面与估值](\#基本面与估值)
\item [资金与交易结构](\#资金与交易结构)
\item [宏观利率背景](\#宏观利率背景)
\item [综合风险与数据局限](\#综合风险与数据局限)
\item [数据与工具说明](\#数据与工具说明)
\item [免责声明](\#免责声明)
\end{itemize}

<a id=''执行摘要''></a>

\subsection{ 执行摘要}
\_本节暂无可用内容。\_

<a id=''核心指标速览''></a>

\subsection{ 核心指标速览}
\begin{table}[htbp]
\centering
\begin{tabular}{lr}
\toprule
指标 & 数值 \\
\midrule
中信一级行业 & 27 \\
最新收盘价 & 39.46 \\
MA20 & 40.5455 \\
MA60 & 39.2218 \\
20 日收益率 & 5.82\% \\
60 日收益率 & -2.29\% \\
RSI14 & 46.1337 \\
20 日均量 & 6908.01 万 \\
PE(TTM) & 54.3554 \\
PB(TTM) & 5.6957 \\
PS(TTM) & 4.0026 \\
股息率(TTM) & 0.46\% \\
总市值 & 616.41 亿 \\
融资余额 & 29.95 亿 \\
融资买入额 & 4.84 亿 \\
\bottomrule
\end{tabular}
\end{table}


<a id=''量价与技术面''></a>

\subsection{ 量价与技术面}

**核心结论**

卧龙电驱（600580.XSHG）股价2026年5月29日收于39.46元，MA20折让2.68\%但站稳MA60上方0.61\%，RSI(14)报46.13处中性偏弱；近20日量增245.74\%领先价涨5.82\%，但近5日量价齐跌7.13\%，短期方向待选择。

\#\#\#\# 图 · RSI 与 MACD

\begin{figure}[htbp]
\centering
\includegraphics[width=0.8\textwidth]{charts/600580_multi_agent_report/technical_indicators.png}
\caption{technical\_indicators.png}
\end{figure}

**图注** RSI 位于中性区间；MACD 位于信号线上方，动能偏强。

\#\#\#\# 表 · 技术指标快照

\begin{table}[htbp]
\centering
\begin{tabular}{ll}
\toprule
指标 & 数值 \\
\midrule
最新收盘价 & 39.46 \\
MA20 & 40.5455 \\
MA60 & 39.2218 \\
20 日收益率 & 5.82\% \\
60 日收益率 & -2.29\% \\
RSI14 & 46.1337 \\
20 日均量 & 6908.01 万 \\
\bottomrule
\end{tabular}
\end{table}


**趋势概览**

\#\#\#\# 图 · 收盘价与 MA20/MA60

\begin{figure}[htbp]
\centering
\includegraphics[width=0.8\textwidth]{charts/600580_multi_agent_report/moving_averages.png}
\caption{moving\_averages.png}
\end{figure}

**图注** 价格位于短均线下、长均线上，呈现震荡整理。

\begin{table}[htbp]
\centering
\begin{tabular}{lll}
\toprule
指标 & 数值（2026-05-29） & 说明 \\
\midrule
最新收盘 & 39.46元 & data.technical.`latest\_close` \\
MA20 & 40.55元 & 收盘价低于MA20 2.68\% \\
MA60 & 39.22元 & 收盘价高于MA60 0.61\% \\
20日收益率 & +5.82\% & 2026-04-29→2026-05-29 \\
60日收益率 & -2.29\% & 2026-03-31→2026-05-29 \\
RSI(14) & 46.13 & 中性偏弱，未超卖 \\
20日均量 & 6,908万股 & data.technical.`avg\_volume\_20d` \\
\bottomrule
\end{tabular}
\end{table}


**近期关键交易日对比**

\begin{table}[htbp]
\centering
\begin{tabular}{llllll}
\toprule
日期 & 收盘价（元） & 涨跌幅 & 成交量（万股） & 成交额（万元） & 换手率（\%） \\
\midrule
2026-05-18 & 40.96 & +1.19\% & 6,748 & 275,537 & 4.33 \\
2026-05-19 & 41.09 & +0.32\% & 5,367 & 220,877 & 3.45 \\
2026-05-20 & 39.62 & -3.58\% & 4,754 & 189,283 & 3.05 \\
2026-05-21 & 40.66 & +2.62\% & 9,786 & 401,132 & 6.28 \\
2026-05-22 & 41.09 & +1.06\% & 5,897 & 240,856 & 3.79 \\
2026-05-25 & 42.49 & +3.41\% & 8,835 & 371,984 & 5.67 \\
2026-05-26 & 44.07 & +3.72\% & 14,907 & 665,437 & 9.57 \\
2026-05-27 & 44.42 & +0.79\% & 11,193 & 495,303 & 7.19 \\
2026-05-28 & 41.68 & -6.17\% & 10,167 & 426,040 & 6.53 \\
2026-05-29 & 39.46 & -5.33\% & 8,051 & 322,026 & 5.17 \\
\bottomrule
\end{tabular}
\end{table}


注：数据源自 data.`price\_recent`（`order\_book\_id`: 600580.XSHG）；换手率来自 data.`turnover\_recent`。

**近期异动**

\begin{itemize}
\item **2026-05-26**：单日涨幅+3.72\%，成交量14,907万股达20日均量2.16倍，换手率9.57\%为20日均换手率（约2.73\%）的3.5倍，成交额66.54亿元较前一日增34\%，属典型放量突破阳线。
\item **2026-05-28**：回落-6.17\%，成交量10,167万股仍处高位，但成交额较5月26日减少36\%至42.60亿元，显示上涨动能衰竭后抛压沉重。
\item **2026-05-29**：续跌-5.33\%，盘中低点38.31元为近20个交易日最低（2026-04-07以来），成交量降至8,051万股，量能收缩但价格加速下行。
\end{itemize}

\#\#\#\# 图 · 回撤曲线

\begin{figure}[htbp]
\centering
\includegraphics[width=0.8\textwidth]{charts/600580_multi_agent_report/drawdown.png}
\caption{drawdown.png}
\end{figure}

**图注** 回撤经历加深后有所修复。

**量价关系**

\#\#\#\# 图 · 收盘价与成交量

\begin{figure}[htbp]
\centering
\includegraphics[width=0.8\textwidth]{charts/600580_multi_agent_report/price_volume.png}
\caption{price\_volume.png}
\end{figure}

**图注** 量价同步走弱。

\begin{itemize}
\item 近20个交易日（2026-04-29至2026-05-29）：成交量从2,329万股增至8,051万股，增幅245.74\%；同期价格从37.29元涨至39.46元，升幅5.82\%。量增领先于价涨，反映资金参与度显著提升（data.`analytical\_evidence`.`price\_windows`.`last\_20d`）。
\item 近5个交易日（2026-05-25至2026-05-29）：价格从42.49元回落至39.46元，跌幅7.13\%；成交量从8,835万股降至8,051万股，降幅8.88\%。量缩与价格下跌方向一致，短期调整中抛压有所减轻。
\end{itemize}

**估值与基本面对照**

\begin{table}[htbp]
\centering
\begin{tabular}{llll}
\toprule
指标 & 2025-09-11 & 2026-05-29 & 变化 \\
\midrule
PE（TTM） & 69.92 & 54.36 & -22.26\% \\
PB（TTM） & 6.44 & 5.70 & -11.49\% \\
净利润增速（TTM） & 191.10\% & 28.48\% & 大幅放缓 \\
毛利率（TTM） & 24.11\% & 25.10\% & +0.99pct \\
资产负债率（\%） & 56.40 & 53.91 & -2.49pct \\
\bottomrule
\end{tabular}
\end{table}


数据源自 data.`factor\_trend`，卧龙电驱（600580.XSHG）隶属电力设备及新能源行业。估值压缩（PE从69.92降至54.36）与基本面边际走弱（净利润增速从191\%降至28.48\%）同步发生，但毛利率微升0.99pct显示主业盈利能力仍稳。

**融资融券观察**

\begin{table}[htbp]
\centering
\begin{tabular}{llll}
\toprule
时点 & 融资余额 & 融资买额 & 说明 \\
\midrule
2025-09-11 & 27.62亿元 & — & 起始观察日 \\
2025-09-18 & — & **峰值15.99亿元** & 历史融资买额峰值 \\
2026-05-28 & 29.95亿元 & 4.84亿元 & 期末融资余额 \\
\bottomrule
\end{tabular}
\end{table}


数据源自 data.`margin\_trajectory`。融资余额自27.62亿元缓步增至29.95亿元（+8.42\%），表明融资资金持续净流入；期末融资买额4.84亿元较峰值15.99亿元大幅回落，融资交易活跃度有所降温。

\begin{itemize}
\item 数据局限：主力资金流向（北向/机构席位）字段为空（data.`capital\_flow`.`row\_count` = 0）；融资融券数据截至2026-05-28，缺少5月29日数据；缺少龙虎榜与大宗交易记录。
\end{itemize}

<a id=''基本面与估值''></a>

\subsection{ 基本面与估值}

**核心结论**

2026年5月29日卧龙电驱（600580.XSHG）市值616.41亿元，PE(TTM)54.36倍较2025年9月11日的69.92倍压缩22.26\%，估值收缩主因净利润增速从191\%回落至28.48\%（低基数效应消退）；同期净利率从5.96\%升至7.42\%（+1.46pct），资产负债率从56.40\%降至53.91\%，盈利能力和财务结构边际改善。

**盈利与成长**

\#\#\#\# 表 · 最新盈利质量因子

\begin{table}[htbp]
\centering
\begin{tabular}{ll}
\toprule
指标 & 最新值 \\
\midrule
毛利率(TTM) & 25.10\% \\
净利率(TTM) & 7.42\% \\
\bottomrule
\end{tabular}
\end{table}


近8个月卧龙电驱归母净利润增长率从195.15\%回落至35.27\%，净利润增长率从191.10\%降至28.48\%，主因2024年同期低基数效应消退。毛利率从24.11\%升至25.10\%（+0.99pct），净利率从5.96\%升至7.42\%（+1.46pct），反映产品结构优化或成本管控见效。毛利润增长率从+30.91\%转为-2.08\%，显示营收增速放缓。营业利润增长率从146.19\%降至24.59\%，增速放缓但仍处双位数区间。

\begin{table}[htbp]
\centering
\begin{tabular}{llll}
\toprule
指标 & 2025-09-11 & 2026-05-29 & 变动 \\
\midrule
毛利率(TTM) & 24.11\% & 25.10\% & +0.99pct \\
净利率(TTM) & 5.96\% & 7.42\% & +1.46pct \\
归母净利润增长率(TTM) & 195.15\% & 35.27\% & -159.88pct \\
净利润增长率(TTM) & 191.10\% & 28.48\% & -162.62pct \\
毛利润增长率(TTM) & 30.91\% & -2.08\% & -32.99pct \\
\bottomrule
\end{tabular}
\end{table}


注：增长率为同比数据，2025年9月对应2024年中报累计口径，2026年5月对应2025年中报口径，基数差异导致增速大幅收敛。

**现金流与勾稽**

应收账款周转率从3.01降至2.91，流动资产周转率从1.14降至1.05，周转效率边际走弱，或因账期延长或收入确认节奏变化。流动比率1.24、速动比率0.94，速动比率低于1暗示存货占比较高情况下短期偿债存在压力。资产负债率从56.40\%降至53.91\%，财务杠杆有所下降。股息率从0.36\%升至0.46\%，2026年4月每股派息0.18元较2025年6月的0.15元增加20\%，反映现金流改善信号。

\#\#\#\# 表 · 最新偿债与流动性

\begin{table}[htbp]
\centering
\begin{tabular}{ll}
\toprule
指标 & 最新值 \\
\midrule
流动比率 & 1.24 倍 \\
速动比率 & 0.94 倍 \\
资产负债率 & 53.91\% \\
\bottomrule
\end{tabular}
\end{table}


注：JSON数据不包含现金流量表（经营现金流/投资现金流/筹资现金流），净利润含金量无法验证，与利润表的勾稽关系无法完整呈现。

**估值水平**

\#\#\#\# 图 · 估值因子走势

\begin{figure}[htbp]
\centering
\includegraphics[width=0.8\textwidth]{charts/600580_multi_agent_report/valuation_factors.png}
\caption{valuation\_factors.png}
\end{figure}

**图注** 多条曲线整体向下，指标同步走弱。

\#\#\#\# 图 · PE/PB 历史分位

\begin{figure}[htbp]
\centering
\includegraphics[width=0.8\textwidth]{charts/600580_multi_agent_report/valuation_percentile.png}
\caption{valuation\_percentile.png}
\end{figure}

PE从69.92压缩至54.36（-22.26\%），PB从6.44降至5.70（-11.49\%），PS维持在4.00附近。市值从655.31亿元降至616.41亿元（-5.94\%），PE收缩幅度（-22.26\%）显著大于市值跌幅（-5.94\%），表明估值压缩是市值下降的主因，而非基本面恶化。

\begin{table}[htbp]
\centering
\begin{tabular}{llll}
\toprule
估值指标 & 2025-09-11 & 2026-05-29 & 变动幅度 \\
\midrule
市值(亿元) & 655.31 & 616.41 & -5.94\% \\
PE(TTM) & 69.92 & 54.36 & -22.26\% \\
PB(TTM) & 6.44 & 5.70 & -11.49\% \\
PS(TTM) & 4.02 & 4.00 & -0.44\% \\
股息率(TTM) & 0.36\% & 0.46\% & +0.10pct \\
\bottomrule
\end{tabular}
\end{table}


**技术面与资金面信号**

RSI(14)为46.13，位于中性偏弱区间。最新收盘价39.46元，较20日均线40.55元折让2.68\%，较60日均线39.22元溢价0.61\%，短期在均线附近获得支撑。融资余额从2025年9月11日的27.62亿元升至2026年5月28日的29.95亿元（+8.44\%），2025年9月18日峰值时杠杆买入达15.99亿元，显示前期高点有大量融资盘被套，近期融资余额上升反映资金关注度回升但或形成抛压。

**数据局限**

\begin{itemize}
\item 仅覆盖2025年9月至2026年5月，无完整年报/季报数据，利润质量无法深度验证
\item 无现金流量表数据，净利润含金量及经营现金流与净利润的背离关系无法判断
\item 无分业务收入结构数据，新能源业务（电力设备及新能源行业归属）占比及边际变化未知
\item 无同行可比公司数据，行业相对估值对比无法进行
\item capital\_flow接口返回行数为0，主力资金净流入/流出数据缺失
\end{itemize}

<a id=''资金与交易结构''></a>

\subsection{ 资金与交易结构}

**核心结论**

卧龙电驱（600580.XSHG）区间融资余额从27.62亿元增至29.95亿元，融资买入峰值出现于2025年9月18日（15.99亿元），当日股价56.40元即为区间最高点，杠杆资金随后被套；2026年5月26日换手率骤升至9.57\%伴随量能放大至1.49亿股，但次日即快速回落，股价由44.07元跌至39.46元，量价走势显示资金集中点火后快速撤离。

\#\#\#\# 图 · 融资融券余额

\begin{figure}[htbp]
\centering
\includegraphics[width=0.8\textwidth]{charts/600580_multi_agent_report/margin_balances.png}
\caption{margin\_balances.png}
\end{figure}

**图注** 融资余额曲线整体上行。

\#\#\#\# 表 · 两融区间变动

\begin{table}[htbp]
\centering
\begin{tabular}{ll}
\toprule
指标 & 数值 \\
\midrule
区间起始 & 2025-09-11 \\
区间结束 & 2026-05-28 \\
融资余额（期初） & 27.62 亿 \\
融资余额（期末） & 29.95 亿 \\
融资余额变动 & 2.33 亿 \\
变动幅度 & 8.42\% \\
融券余额（期初） & 0.37 亿 \\
融券余额（期末） & 0.28 亿 \\
融资买入峰值 & 15.99 亿（2025-09-18） \\
\bottomrule
\end{tabular}
\end{table}


**成交与活跃度**

\#\#\#\# 表 · 成交活跃度

\begin{table}[htbp]
\centering
\begin{tabular}{ll}
\toprule
指标 & 数值 \\
\midrule
20 日均量 & 6908.01 万 \\
近12日日均成交额 & 34.61 亿 \\
近12日最高成交额 & 66.54 亿（2026-05-26） \\
近12日最低成交额 & 18.78 亿（2026-05-14） \\
近12日换手率区间 & 3.03\% \textasciitilde{} 9.57\% \\
主力资金流向 & 数据缺失 \\
\bottomrule
\end{tabular}
\end{table}


区间观察窗口为2025-09-11至2026-05-29（共169个交易日）。20日区间日均成交量69,080,123股（约6,907万股），对应 `analytical\_evidence.`price\_snapshot`.avg\_volume\_20d`。

近5日（2026-05-25至2026-05-29）量价齐缩：收盘价由42.49元跌至39.46元（-7.13\%），成交量由88,351,702股降至80,506,725股（-8.88\%），成交额由37.20亿元降至32.20亿元（-13.43\%），数据来源于 逐日追溯。

\begin{table}[htbp]
\centering
\begin{tabular}{lllll}
\toprule
日期 & 收盘价（元） & 成交量（股） & 成交额（元） & 当日换手率（\%） \\
\midrule
2026-05-25 & 42.49 & 88,351,702 & 3,719,841,932 & 5.67 \\
2026-05-26 & 44.07 & 149,073,340 & 6,654,366,401 & 9.57 \\
2026-05-27 & 44.42 & 111,925,572 & 4,953,032,638 & 7.19 \\
2026-05-28 & 41.68 & 101,673,802 & 4,260,403,847 & 6.53 \\
2026-05-29 & 39.46 & 80,506,725 & 3,220,255,872 & 5.17 \\
\bottomrule
\end{tabular}
\end{table}


2026年5月26日出现显著异常：成交量环比前日放大68.7\%至1.49亿股，成交额达66.54亿元，当日换手率9.57\%为区间峰值，价格触及46.49元（次高点）。次日（5-27）量能仍维持1.12亿股高位后，5月28日起连续回落，5-29当日跌幅5.33\%收于39.46元。该走势符合''资金集中点火后快速撤离''特征，来源于 `turnover\_recent.（9.5705\%）与 `price\_change\_rate\_recent` 逐日对比。

**融资融券**

\#\#\#\# 图 · 融资余额与买卖额

\begin{figure}[htbp]
\centering
\includegraphics[width=0.8\textwidth]{charts/600580_multi_agent_report/margin_enhanced.png}
\caption{margin\_enhanced.png}
\end{figure}

\#\#\#\# 表 · 融资融券快照

\begin{table}[htbp]
\centering
\begin{tabular}{ll}
\toprule
指标 & 数值 \\
\midrule
统计日期 & 2026-05-28 \\
融资余额 & 29.95 亿 \\
融券余额 & 0.28 亿 \\
融资买入额 & 4.84 亿 \\
融资偿还额 & 4.68 亿 \\
融券余量 & 67.51 万 \\
两融余额合计 & 30.23 亿 \\
\bottomrule
\end{tabular}
\end{table}


融资余额区间（2025-09-11至2026-05-28）从27.62亿元升至29.95亿元，净增2.33亿元，来源于 `margin\_trajectory.margin\_balance\_start` 与 `margin\_trajectory.margin\_balance\_end`。

2025年9月18日融资买入峰值15.99亿元（`margin\_trajectory.peak\_buy\_on\_margin\_value`），当日 `period\_extremes.high.close` 为56.40元即为区间最高点，杠杆资金随后被套。

近两周融资余额持续攀升，来源于 `securities\_margin\_recent` 逐日追溯：

\begin{table}[htbp]
\centering
\begin{tabular}{llll}
\toprule
日期 & 融资余额（元） & 融资买入额（元） & 融券余额（元） \\
\midrule
2026-05-13 & 2,493,345,867 & 170,603,222 & 19,240,642 \\
2026-05-20 & 2,576,477,441 & 203,350,643 & 20,259,330 \\
2026-05-26 & 2,880,378,606 & 791,947,480 & 31,790,379 \\
2026-05-28 & 2,994,506,304 & 483,536,809 & 28,136,542 \\
\bottomrule
\end{tabular}
\end{table}


2026-05-13至2026-05-28期间，融资余额由24.93亿元升至29.95亿元（+19.9\%）；2026-05-26融资买入额达7.92亿元为近期单日最高，当日股价44.07元，次日（5-27）融资买入仍维持5.89亿元高位后，股价随后下跌。融券余额从2026-05-13的1,924万元降至2026-05-28的2,814万元（5-26当天融券余额曾升至3,179万元），方向不统一，无法确认明确做空压力。

融资余额攀升但股价未创新高，显示杠杆资金支撑与价格走势出现背离，原因可能为套牢盘补仓或资金博弈短期反弹。

**股东与股本**

\#\#\#\# 表 · 股本结构快照

\begin{table}[htbp]
\centering
\begin{tabular}{ll}
\toprule
项目 & 数值 \\
\midrule
统计日期 & 2026-05-29 \\
总股本 & 15.62 亿 \\
流通 A 股 & 15.58 亿 \\
自由流通股本 & 7.95 亿 \\
自由流通占比 & 50.9\% \\
非自由流通占比 & 49.1\% \\
\bottomrule
\end{tabular}
\end{table}


\begin{table}[htbp]
\centering
\begin{tabular}{lll}
\toprule
类别 & 股数（股） & 占比 \\
\midrule
总股本 & 1,562,117,511 & 100\% \\
A股流通 & 1,557,632,511 & 99.71\% \\
非流通A股 & 4,485,000 & 0.29\% \\
自由流通股本 & 795,154,658 & 50.90\% \\
\bottomrule
\end{tabular}
\end{table}


区间内股本结构保持稳定，无送转记录，来源于 推算。

**分红与资金成本**

近三次分红（含税）呈逐年递增趋势，来源于 `factor\_snapshot.dividend\_yield\_ttm` 及 `factor\_trend` 逐期对比：

\begin{table}[htbp]
\centering
\begin{tabular}{lll}
\toprule
分红季度 & 除权除息日 & 每10股派息（元） \\
\midrule
`2024q4` & 2025-06-27 & 1.5 \\
`2025q4` & 2026-04-30 & 1.8 \\
\bottomrule
\end{tabular}
\end{table}


以最新收盘价39.46元计算，TTM股息率约0.456\%（`factor\_snapshot.dividend\_yield\_ttm`），对融资利息成本（约5\%-6\%）覆盖能力有限（0.456\% ÷ 5\% ≈ 9.1\%），说明股息收益无法覆盖杠杆资金利息支出，融资持有卧龙电驱存在负carry风险。

**数据局限**

\begin{itemize}
\item 融资数据未提供
\end{itemize}

<a id=''宏观利率背景''></a>

\subsection{ 宏观利率背景}

**核心结论**

短端利率隔夜品种Shibor ON在2026年5月29日为1.324\%，较20日前（2026-04-29）上行8.1bp至1.324\%；1Y Shibor下降0.9bp至1.466\%；国债1Y/10Y/30Y收益率分别为1.158\%/1.709\%/2.205\%，期限利差（1Y-10Y）约55bp。

\#\#\#\# 图 · Shibor 利率

\begin{figure}[htbp]
\centering
\includegraphics[width=0.8\textwidth]{charts/600580_multi_agent_report/shibor_rates.png}
\caption{shibor\_rates.png}
\end{figure}

**图注** 短端利率整体下行，波动相对平稳。

**短端利率：隔夜品种小幅回升**

\begin{table}[htbp]
\centering
\begin{tabular}{lllll}
\toprule
Shibor期限 & 2026-04-29 & 2026-05-29 & 变化（bp） & 来源字段 \\
\midrule
ON（隔夜） & 1.243\% & 1.324\% & +8.1 & `macro\_snapshots`.`interbank\_20d\_ago`.ON / `interbank\_latest`.ON \\
1W & 1.373\% & 1.387\% & +1.4 & `macro\_snapshots`.`interbank\_20d\_ago`.1W / `interbank\_latest`.1W \\
2W & 1.376\% & 1.387\% & +1.1 & `macro\_snapshots`.`interbank\_20d\_ago`.2W / `interbank\_latest`.2W \\
1M & 1.395\% & 1.390\% & -0.5 & `macro\_snapshots`.`interbank\_20d\_ago`.1M / `interbank\_latest`.1M \\
3M & 1.420\% & 1.408\% & -1.2 & `macro\_snapshots`.`interbank\_20d\_ago`.3M / `interbank\_latest`.3M \\
6M & 1.447\% & 1.429\% & -1.8 & `macro\_snapshots`.`interbank\_20d\_ago`.6M / `interbank\_latest`.6M \\
9M & 1.460\% & 1.450\% & -1.0 & `macro\_snapshots`.`interbank\_20d\_ago`.9M / `interbank\_latest`.9M \\
1Y & 1.475\% & 1.466\% & -0.9 & `macro\_snapshots`.`interbank\_20d\_ago`.1Y / `interbank\_latest`.1Y \\
\bottomrule
\end{tabular}
\end{table}


**图表解读：** 隔夜品种（ON）上行8.1bp反映短期流动性边际收紧；1M-1Y品种全线小幅回落，显示中期资金面仍相对宽松。两者走势背离，表明市场对短期流动性波动的预期与中期利率趋势存在分歧。

**收益率曲线结构**

\#\#\#\# 图 · 国债收益率走势

\begin{figure}[htbp]
\centering
\includegraphics[width=0.8\textwidth]{charts/600580_multi_agent_report/gov_yield_trend.png}
\caption{gov\_yield\_trend.png}
\end{figure}

**图注** 多条曲线整体向下，指标同步走弱。

\begin{table}[htbp]
\centering
\begin{tabular}{lll}
\toprule
期限 & 2026-05-29收益率 & 来源字段 \\
\midrule
隔夜（0S） & 0.98\% & `yield\_curve\_latest`.0S \\
1M & 1.043\% & `yield\_curve\_latest`.1M \\
3M & 1.061\% & `yield\_curve\_latest`.3M \\
6M & 1.098\% & `yield\_curve\_latest`.6M \\
1Y & 1.158\% & `yield\_curve\_latest`.1Y \\
2Y & 1.226\% & `yield\_curve\_latest`.2Y \\
3Y & 1.271\% & `yield\_curve\_latest`.3Y \\
5Y & 1.419\% & `yield\_curve\_latest`.5Y \\
7Y & 1.558\% & `yield\_curve\_latest`.7Y \\
10Y & 1.709\% & `yield\_curve\_latest`.10Y \\
15Y & 2.017\% & `yield\_curve\_latest`.15Y \\
20Y & 2.203\% & `yield\_curve\_latest`.20Y \\
30Y & 2.205\% & `yield\_curve\_latest`.30Y \\
\bottomrule
\end{tabular}
\end{table}


**图表解读：** 1Y-10Y期限利差约55bp（10Y 1.709\% - 1Y 1.158\%），曲线形态陡峭。30Y国债收益率（2.205\%）与20Y（2.203\%）几乎持平，长期限品种出现平坦化特征，反映市场对长期增长预期的分歧。

**卧龙电驱（600580.XSHG）估值与利率环境的逻辑联系**

\begin{table}[htbp]
\centering
\begin{tabular}{llll}
\toprule
估值指标 & 数值 & 来源字段 & 与债券收益率关联 \\
\midrule
股息率TTM & 0.456\% & factor.`dividend\_yield\_ttm` & 10Y国债1.709\%对卧龙电驱股息率（0.456\%）形成约125bp溢价，股票现金回报吸引力相对有限 \\
PE（TTM） & 54.36x & factor.`pe\_ratio\_ttm` & 无直接逻辑联系 \\
PB（TTM） & 5.70x & factor.`pb\_ratio\_ttm` & 无直接逻辑联系 \\
\bottomrule
\end{tabular}
\end{table}


\begin{itemize}
\item 卧龙电驱隶属电力设备及新能源行业（`first\_industry\_code`: ''27''），行业资产偏重，利率变化通过融资成本影响资本开支计划。
\end{itemize}

**融资余额变化与利率环境对照**

\begin{table}[htbp]
\centering
\begin{tabular}{lll}
\toprule
时点 & 融资余额（亿元） & 来源字段 \\
\midrule
2025-09-11 & 27.62 & `margin\_trajectory`.`margin\_balance\_start` \\
2026-05-28 & 29.95 & `margin\_trajectory`.`margin\_balance\_end` \\
\bottomrule
\end{tabular}
\end{table}


\begin{itemize}
\item 融资余额从27.62亿元增至29.95亿元，显示杠杆资金持续净流入；但与利率环境变化的逻辑联系无法从现有数据直接建立，属于数据局限。
\end{itemize}

**数据局限**

\begin{itemize}
\item ，主力资金流向无法追踪
\end{itemize}

<a id=''综合风险与数据局限''></a>

\subsection{ 综合风险与数据局限}

\_本节暂无可用内容。\_

<a id=''数据与工具说明''></a>

\subsection{ 数据与工具说明}
\begin{itemize}
\item 数据区间：2025-09-11 至 2026-05-29
\item 计划使用的米筐函数：get\_price, get\_price\_change\_rate, get\_turnover\_rate, get\_capital\_flow, get\_factor, get\_securities\_margin, get\_dividend, get\_shares, get\_instrument\_industry, is\_suspended, is\_st\_stock, get\_interbank\_offered\_rate, get\_yield\_curve, get\_pit\_financials\_ex
\item 数据质量：9 项可用；缺失 基准指数, 资金流向, 大宗交易, 年报三表
\item 生成时间：2026-05-31 22:40
\end{itemize}

<a id=''免责声明''></a>

\subsection{ 免责声明}
本报告仅供课程研究与信息展示，不构成投资建议。\end{document}
